\chapter{Race conditions}
\label{ch:races}

Two threads that only read shared data never disturb each other. As soon
as one of them writes, the result can depend on the exact order in which
the threads' instructions happen to interleave --- an order nobody
controls. This chapter explains why even \code{counter++} is not safe,
how to reason about interleavings, and which guarantees (``happens
before'') you can actually rely on.

\section{The lost update}

\code{counter++} looks like one step. The CPU executes three:

\begin{lstlisting}
mov  eax, [counter]     ; load
add  eax, 1             ; modify (in a register - private to the thread)
mov  [counter], eax     ; store
\end{lstlisting}

If thread A is interrupted (or simply runs on another core) between its
load and its store, both threads can load the same old value:

\begin{center}
\begin{tabular}{llc}
\toprule
\textbf{Thread A} & \textbf{Thread B} & \textbf{counter} \\
\midrule
load (reads 5)    &                    & 5 \\
                  & load (reads 5)     & 5 \\
                  & add, store 6       & 6 \\
add, store 6      &                    & \textbf{6} \\
\bottomrule
\end{tabular}
\end{center}

Two increments, one effect: a \emph{lost update}. Two threads with three
instructions each can interleave in $\binom{6}{3} = 20$ ways; only the 2
in which one thread's three instructions all come before the other's are
correct (\module{03}'s \code{interleave.c} lists all 20). With a million
increments per thread, the final count is usually far too small --- and
different in every run.

\section{Check-then-act}

A race needs no \code{++}. Any sequence ``look at shared state, then act
on what you saw'' breaks if the state changes in between:

\begin{lstlisting}
if (account->balance >= amount)        /* check */
  account->balance -= amount;          /* act: another thread may have
                                          withdrawn in between */
\end{lstlisting}

The same pattern appears as ``if the file does not exist, create it'',
``if the cache has no entry, compute and insert it'', ``if the pointer is
NULL, allocate''. In a kernel the dangerous variant is
\emph{time-of-check to time-of-use} (TOCTOU): the kernel checks a value in
user memory, then reads it again to use it --- and another user thread
changed it in between (chapter~7).

\section{Critical sections and atomicity}

A \emph{critical section} is a piece of code that accesses shared state
and must not run concurrently with other critical sections on the same
state. Making it \emph{atomic} means: every other thread sees either none
of its effects or all of them, never an intermediate state. There are
three ways to get there (\module{03}'s \code{three\_fixes.c} compares
them):

\begin{enumerate}
  \item \textbf{Lock it}: a mutex around the whole check-and-act
    (chapter~4).
  \item \textbf{Use an atomic operation}: \code{atomic\_fetch\_add} does
    load--add--store as one indivisible instruction (\code{lock xadd}).
    Only works when the whole critical section is a single such operation.
  \item \textbf{Do not share}: give each thread its own copy and combine
    afterwards (chapter~2's partition-then-combine).
\end{enumerate}

Invariants are the key to finding critical sections: ``the sum of all
balances is constant'', ``\code{count} equals the number of elements in
the list''. Every code path that temporarily breaks an invariant is a
critical section.

\section{Data race vs.\ race condition}

The two words are often mixed up:

\begin{description}
  \item[Data race] (a precise, mechanical definition): two threads access
    the same memory location, at least one access is a write, and the two
    accesses are not ordered by happens-before (no lock, no atomic, no
    join between them). In C and C++ a data race is \emph{undefined
    behaviour} --- the compiler may assume it never happens.
    ThreadSanitizer finds exactly these.
  \item[Race condition] (a semantic property): the program's correctness
    depends on timing. A program can have a race condition without a data
    race (every access is locked, but the lock is released between check
    and act) and a data race without visible harm (a statistics counter
    that may be a little off --- still undefined behaviour).
\end{description}

\section{Happens-before}

Which writes of thread A is thread B guaranteed to see? Only those that
\emph{happen before} B's read, where happens-before is built from:

\begin{itemize}
  \item program order within one thread;
  \item unlocking a mutex happens before the next lock of the same mutex;
  \item \code{pthread\_create} happens before the first instruction of the
    new thread; the thread's last instruction happens before
    \code{pthread\_join} returns;
  \item a release store to an atomic happens before an acquire load that
    reads it (chapter~4);
  \item transitivity: A before B and B before C gives A before C.
\end{itemize}

Anything else is unordered, and without an ordering there is no
guarantee at all --- not even that B ever sees A's write. The compiler
may keep a variable in a register for the whole loop:

\begin{lstlisting}
while (!done) {}       /* with a plain `int done`, -O2 may turn this into
                          `if (!done) for (;;);` - it never re-reads it */
\end{lstlisting}

and the CPU may make stores visible to other cores in a different order
than the program wrote them (chapter~4). \code{volatile} stops the
compiler from caching, but orders nothing between threads; use atomics or
locks.

\section{Why testing is not enough}

A race that loses one update in a million runs passes every normal test.
It also depends on the machine: more cores, a different load, an
optimised build or an extra \code{printf} change the timing. Hence the
modules test twice --- under ThreadSanitizer, which reports the data race
whenever the unordered accesses \emph{happen}, whether or not they
collide in that run, and optimised at full speed, where broken invariants
become visible.

\begin{swebbox}
SWEB is preemptive: the timer interrupt can switch threads at almost any
instruction, even in kernel code (\code{syscallHandler} enables
interrupts before calling \code{Syscall::syscallException}). After A1,
several threads of one process can be in the kernel at the same time,
touching the same process structures --- the thread list, the address
space, the file descriptors. Every such structure needs a lock, and every
``check, then act'' in your \code{pthread\_join} (``is the thread
finished? if not, sleep'') is a critical section. The A1 plan says it
bluntly: \emph{races appear on the tenth run, not the first.}
\end{swebbox}

\begin{keybox}
\begin{itemize}
  \item \code{x++} is load--modify--store; interleavings lose updates.
  \item Check-then-act is a race even without arithmetic; in the kernel
    it is TOCTOU.
  \item Find critical sections through invariants; fix them by locking,
    atomics, or not sharing.
  \item Data race = unordered conflicting accesses (undefined behaviour);
    race condition = timing-dependent correctness.
  \item Only happens-before edges (locks, create/join, atomics) guarantee
    visibility; \code{volatile} is not synchronisation.
\end{itemize}
\end{keybox}
